\chapter{Sześć operacji snopowych, dualność Verdiera i cykle specjalizacji}
\label{ch:xvi-szesc-operacji}

\section{Motywacja: dlaczego zwykły obraz bezpośredni nie wystarcza?}
W Rozdziale~\ref{ch:snopy-verdier-v12} skonstruowaliśmy snopy,
przedłużenie przez zero i kohomologię z nośnikami zwartymi, a w
Rozdziale~\ref{ch:xiv-hodge-monodromia} napotkaliśmy cykle pobliskie.
Potrzebujemy teraz języka, w którym operacje te tworzą jeden spójny
rachunek. Problem jest geometryczny: jeśli $f:X\to Y$ nie jest właściwe,
przesuwanie informacji z $X$ na $Y$ może zachować sekcje rozciągające się
do nieskończoności albo odrzucić te, których nośnik nie jest właściwy
nad $Y$. Różnicę kodują odpowiednio $Rf_*$ i $Rf_!$.

Ustalamy ciało $k$ i korzystamy z kompleksów snopów $k$-przestrzeni
wektorowych, numerowanych kohomologicznie. Dla przestrzeni $X$
piszemy $D(X)=D(\operatorname{Sh}(X;k))$ i $D^b_c(X)$ dla
ograniczonych kompleksów o konstruktybilnych snopach kohomologii.
W twierdzeniach konstruktybilnych $X$ i $Y$ oznaczają
analityfikacje zespolonych odmian algebraicznych o skończonej
stratyfikacji algebraicznej, a $f:X\\to Y$ jest analityfikacją
morfizmu algebraicznego. Na takich przestrzeniach posługujemy się
kategorią $D^b_c(X;k)$ kompleksów o konstruktybilnych snopach
kohomologii. W przykładach topologicznych można szerzej dopuścić
dobre, lokalnie zwarte przestrzenie Hausdorffa o skończonym wymiarze
kohomologicznym, pod warunkiem sprawdzenia właściwych hipotez
konkretnego twierdzenia. Kategoria $D(X)$ wszystkich kompleksów nie
musi zachowywać ograniczoności ani konstruktybilności przy dowolnej
mapie $f$. Twierdzenia formalizmu sześciu operacji przyjmujemy w
sprecyzowanej kategorii, a dowody przypadków szczególnych zaznaczamy.
Dla kompleksu $A$ piszemy $\\mathcal H^q(Rf_*A)$ na oznaczenie
jego $q$-tego snopu kohomologii; symbol $R^qf_*\\mathcal F$
stosujemy do snopu $\\mathcal F$ skoncentrowanego w stopniu zerowym. Współczynniki w ciele upraszczają
iloczyn tensorowy, lecz nie eliminują konieczności pochodnego
wewnętrznego Homu.

\begin{definicja}[Przesunięcie i stożek]
Dla kompleksu kohomologicznego $K$ definiujemy $K[1]^n=K^{n+1}$
i $d_{K[1]}=-d_K$. Jeśli $u:A\to B$ jest morfizmem kompleksów,
jego stożek $\operatorname{Cone}(u)$ ma
$\operatorname{Cone}(u)^n=B^n\oplus A^{n+1}$ i
$d(b,a)=(d_Bb+u(a),-d_Aa)$. W kategorii pochodnej otrzymujemy
trójkąt wyróżniony
$A\xrightarrow{u}B\to\operatorname{Cone}(u)\to A[1]$.
\end{definicja}
\begin{lemat}[Długi ciąg kohomologii stożka]
Dla każdego morfizmu kompleksów $u:A\to B$ istnieje naturalny
ciąg dokładny
$\cdots\to H^n(A)\to H^n(B)\to
H^n(\operatorname{Cone}(u))\to H^{n+1}(A)\to\cdots$.
\end{lemat}
\begin{proof}
Rzut $\operatorname{Cone}(u)^n\to A^{n+1}$ daje krótki ciąg
dokładny kompleksów $0\to B\to\operatorname{Cone}(u)\to A[1]\to0$.
Morfizm łączący na kohomologii jest indukowany przez $u$ (z
konwencją znaku zależną od przyjętej definicji stożka). Standardowa
konstrukcja długiego ciągu dokładnego daje tezę.
\end{proof}

\section{Sześć funktorów i ich adjunkcje}
\begin{definicja}[Dwie pary funktorów związanych z odwzorowaniem]
Niech $f:X\to Y$ będzie mapą w rozważanej kategorii przestrzeni.
Przez $f^*=f^{-1}$ oznaczamy dokładny funktor obrazu odwrotnego
snopów $k$-przestrzeni wektorowych, a przez $Rf_*$ jego prawy
pochodny obraz bezpośredni. Dla snopu $\mathcal F$ obraz bezpośredni
jest dany przez $(f_*\mathcal F)(V)=\mathcal F(f^{-1}V)$.
Przez $f_!\mathcal F$ oznaczamy snop sekcji o nośniku właściwym
nad otwartym podzbiorem bazy: na $V\subseteq Y$ rozważamy
sekcje nad $f^{-1}V$, których nośnik odwzorowuje się właściwie na $V$,
i dokonujemy snopifikacji, jeśli jest potrzebna. Prawy funktor
pochodny oznaczamy $Rf_!$. W tych samych założeniach $f^!$
jest prawym adjunktem do $Rf_!$ na odpowiednich kategoriach
pochodnych.
\end{definicja}
\begin{uwaga}[Cztery operacje geometryczne]
Zapis $f^*$ nie jest tutaj rozszerzaniem skalarów po pierścieniach
funkcji; pracujemy ze snopami $k$-przestrzeni wektorowych.
Dla $f$ właściwego zachodzi $Rf_!\simeq Rf_*$. Zwykły
$f_*$ jest lewostronnie dokładny, zaś $f_!$ i $Rf_!$
nie muszą zachowywać krótkich ciągów dokładnych bez przejścia
na kategorię pochodną.
\end{uwaga}
\begin{definicja}[Tensor i wewnętrzny Hom]
Dla $A,B\in D(X)$ przez $A\otimes_k^L B$ oznaczamy
pochodny iloczyn tensorowy snopów, a przez
$R\mathcal Hom_X(A,B)$ pochodny \emph{wewnętrzny Hom},
charakteryzowany przez adjunkcję
\[
\operatorname{Hom}_{D(X)}(A\otimes_k^L B,C)
\cong\operatorname{Hom}_{D(X)}
\bigl(A,R\mathcal Hom_X(B,C)\bigr).
\]
W naszym przypadku każdy $k$-moduł jest płaski, toteż
$\otimes_k^L$ można obliczać zwykłym tensorem kompleksów.
Natomiast $R\mathcal Hom_X$ jest funktorem pochodnym w kategorii
snopów: zwykły wewnętrzny Hom nie wystarcza.
\end{definicja}
\begin{twierdzenie}[Podstawowe adjunkcje: sformułowanie formalizmu]
\label{thm:xvi-adjunctions}
Niech $f:X\to Y$ będzie morfizmem, dla którego obowiązuje
formalizm sześciu operacji. Dla $A\in D(X)$ i $B\in D(Y)$
istnieją naturalne izomorfizmy
\[
\operatorname{Hom}_{D(X)}(f^*B,A)
\cong\operatorname{Hom}_{D(Y)}(B,Rf_*A),
\quad
\operatorname{Hom}_{D(Y)}(Rf_!A,B)
\cong\operatorname{Hom}_{D(X)}(A,f^!B).
\]
Razem z $\otimes^L$ i $R\mathcal Hom$ tworzą one sześć
operacji $f^*,Rf_*,Rf_!,f^!,\otimes^L,R\mathcal Hom$.
\end{twierdzenie}
\begin{proof}[Część dowiedziona i część przyjęta]
Pierwsza adjunkcja w stopniu zerowym wynika z definicji obrazu
odwrotnego: morfizm $f^{-1}\mathcal G\to\mathcal F$
odpowiada morfizmowi $\mathcal G\to f_*\mathcal F$.
Ponieważ $f^{-1}$ jest dokładny, adjunkcja przechodzi na
kompleksy przez zastosowanie rezolwent injektywnych do
argumentu $\mathcal F$.
Druga adjunkcja \emph{definiuje} szczególny obraz odwrotny $f^!$;
świadomego dowodu jego istnienia dla ogólnych map nie można
zastąpić manipulacją symbolami. Istnienie i własności $f^!$
w przyjętej kategorii stanowią tu standardowe twierdzenie
formalizmu dualności Verdiera.
\end{proof}

\section{Otwarte i zamknięte włożenia: jawne obliczenia}
Niech $i:Z\hookrightarrow X$ będzie włożeniem zamkniętym,
a $j:U=X\setminus Z\hookrightarrow X$ otwartym.
Dla snopu $\mathcal F$ na $X$ definiujemy snop sekcji
z nośnikiem w $Z$ przez
$(\underline\Gamma_Z\mathcal F)(V)=
\{s\in\mathcal F(V):\operatorname{supp}(s)\subseteq Z\cap V\}$
na otwartych $V\subseteq X$. Odróżniamy go od przestrzeni
sekcji globalnych $\Gamma_Z(X,\mathcal F)$.
Wówczas dla $A\in D(X)$ zachodzi
$i_*i^!A\simeq R\underline\Gamma_Z(A)$;
jest to pochodny kompleks snopów z nośnikiem w $Z$.

\begin{lemat}[Pierwszy ciąg lokalizacyjny na poziomie snopów]
Dla każdego snopu $\mathcal F$ na $X$ mamy dokładny ciąg
\[
0\longrightarrow j_!j^*\mathcal F
\longrightarrow\mathcal F\longrightarrow i_*i^*\mathcal F
\longrightarrow0.
\]
\end{lemat}
\begin{proof}
Wystarczy sprawdzić dokładność na kiełkach.
Dla $x\in U$ kiełki trzech pierwszych wyrazów to
$\mathcal F_x\to\mathcal F_x\to0$.
Dla $x\in Z$ są to $0\to\mathcal F_x\to\mathcal F_x$.
W obu przypadkach otrzymujemy ciąg dokładny.
\end{proof}
\begin{twierdzenie}[Dwa trójkąty lokalizacyjne]
\label{thm:xvi-localization}
Dla $A\in D^b_c(X)$ zachodzą naturalne trójkąty wyróżnione
\[
 j_!j^*A\to A\to i_*i^*A\to (j_!j^*A)[1],
\]
\[
 i_*i^!A\to A\to Rj_*j^*A\to (i_*i^!A)[1].
\]
\end{twierdzenie}
\begin{proof}[Pierwszy trójkąt i status drugiego]
Pierwszy trójkąt pochodzi z poprzedniego krótkiego ciągu dokładnego
zastosowanego składnikowo do kompleksu; funktory $j_!,j^*,i_*,i^*$
są tutaj dokładne. Dla drugiego wybierzmy rezolwentę
injektywną $A\to I^\bullet$. Snopy injektywne są wiotkie
(\emph{flabby}), czyli ograniczenie sekcji do dowolnego
mniejszego zbioru otwartego jest surjekcją.
Dla każdego otwartego $V\subseteq X$ mamy zatem krótki ciąg
$0\to\underline\Gamma_Z(I^q)(V)
\to I^q(V)\to I^q(V\cap U)\to0$.
Po zebraniu stopni i przejściu na snopy otrzymujemy
krótki ciąg kompleksów
$0\to\underline\Gamma_Z(I^\bullet)
\to I^\bullet\to j_*j^*I^\bullet\to0$.
Ponieważ $j^*$ ma dokładny lewy adjunkt $j_!$,
zachowuje snopy injektywne; dlatego
$j_*j^*I^\bullet$ oblicza $Rj_*j^*A$.
Pierwszy kompleks jest modelem dla
$R\underline\Gamma_Z(A)\simeq i_*i^!A$.
Powstały trójkąt ma zatem żądaną postać.
Korzystaliśmy tutaj z istnienia rezolwent injektywnych
oraz identyfikacji pochodnych sekcji z nośnikiem
z $i_*i^!$, której konstrukcja wynika z adjunkcji.
\end{proof}
\begin{przyklad}[Zwykły i wyjątkowy obraz punktu na prostej]
\label{ex:xvi-point-line}
Niech $X=\R$, $Z=\{0\}$, $U=\R\setminus\{0\}$,
$A=k_{\R}$. Wówczas $i^*A\simeq k_Z$, ale
$i^!A\simeq k_Z[-1]$. Aby otrzymać drugą równość,
obliczamy lokalną kohomologię:
$H^q(\R,\R\setminus\{0\};k)$ jest równa $k$ dla $q=1$
i zerowa w pozostałych stopniach. Rzeczywiście, długi ciąg
pary i $H^0(\R\setminus\{0\};k)\cong k\oplus k$
pokazują, że $H^1$ jest ilorazem $k^2/\Delta k\cong k$.
Daje to $i^!k_{\R}\simeq k_Z[-1]$.
\end{przyklad}
\begin{przyklad}[Dlaczego $j_!\ne j_*$?]
Dla $j:(0,1)\hookrightarrow\R$ i $A=k_{(0,1)}$
kiełek $(j_!A)_0$ jest zerowy, natomiast
$(j_*A)_0\cong k$. Z drugiego trójkąta wynika również,
że $Rj_*$ może mieć niezerowe wyższe snopy kohomologii,
jeśli lokalne dopełnienie jest topologicznie nietrywialne.
W tym konkretnym przykładzie wyższe $R^qj_*A$ znikają.
\end{przyklad}


\section{Jak rozpoznać $f^!$? Trzy podstawowe przypadki}
\begin{stwierdzenie}[Włożenie otwarte]
\label{prop:xvi-open-shriek}
Dla włożenia otwartego $j:U\hookrightarrow X$ zachodzi
$j^!\simeq j^*$ na kategorii pochodnej snopów $k$-przestrzeni
wektorowych.
\end{stwierdzenie}
\begin{proof}
Dla snopu $A$ na $U$ i snopu $B$ na $X$ naturalne
przyporządkowanie morfizmów daje
$\operatorname{Hom}_X(j_!A,B)\cong
\operatorname{Hom}_U(A,j^*B)$.
Przyporządkowanie jest zgodne ze złożeniami i wynika
z definicji sekcji przedłużonych przez zero.
Ponieważ $j_!$ i $j^*$ są dokładne (co można sprawdzić
na kiełkach), ta sama adjunkcja obowiązuje dla kompleksów.
Natomiast $j^!$ jest z definicji prawym adjunktem do $Rj_!=j_!$.
Jednoznaczność adjunktu daje $j^!\simeq j^*$.
\end{proof}

\begin{twierdzenie}[Gładka czystość: wersja analityczno-topologiczna]
\label{thm:xvi-smooth-purity}
Niech $f:M\to N$ będzie gładką submersją rzeczywistych
rozmaitości, której włókna mają stały wymiar $d$, i niech
$\operatorname{or}_{M/N}$ będzie systemem ich orientacji.
Dla konstruktybilnego kompleksu $B\in D^b_c(N)$ zachodzi
naturalny izomorfizm
\[
 f^!B\simeq f^*B\otimes_k\operatorname{or}_{M/N}[d].
\]
Dla orientowanej submersji system orientacji jest trywialny;
w szczególności dla $a:\R^d\to\{*\}$ mamy
$a^!k\simeq k_{\R^d}[d]$.
\end{twierdzenie}
\begin{proof}[Rachunek lokalny i zaznaczenie zależności]
Z twierdzenia o lokalnej postaci submersji dostajemy
lokalne współrzędne, w których $f$ jest projekcją
$V\times\R^d\to V$. Zorientowane włókno posiada lokalną
klasę fundamentalną w
$H_c^d(\R^d;k)\cong k$, więc przesunięcie o $d$
i skręcenie systemem orientacji są konieczne, aby
adjunkcja $Rf_!\dashv f^!$ zachodziła w modelu
produktowym. Na przecięciach kart zmiana lokalnej
orientacji mnoży generator przez $+1$ lub $-1$,
co składa się na $\operatorname{or}_{M/N}$.
Zgodność lokalnych identyfikacji z ogólną adjunkcją i
ich kanoniczne sklejenie są częścią twierdzenia o
czystości (purity) w dualności Verdiera. Nie wywodzimy
tej globalnej części wyłącznie z lokalnego rachunku grup.
\end{proof}
\begin{przyklad}[Ta sama przestrzeń, różne obrazy do punktu]
Dla $a:\R\to\{*\}$ mamy
$Ra_*k_{\R}\simeq k$ oraz $Ra_!k_{\R}\simeq k[-1]$.
Pierwsza to zwykła kohomologia $\R$, druga to
kohomologia z nośnikami zwartymi; jest ona skoncentrowana
w stopniu $1$. Ponadto $a^!k\simeq k_{\R}[1]$.
Te trzy wzory pokazują jednocześnie, dlaczego
$f_*$ nie zastępuje $f_!$ i dlaczego $f^!$
nie powinno się utożsamiać z $f^*$.
\end{przyklad}

\section{Kohomologia z nośnikiem zwartym i obiekt dualizujący}
\begin{definicja}[Kompleks dualizujący]
Dla $a:X\to\{*\}$ definiujemy $\omega_X=a^!k$,
a dla $A\in D^b_c(X)$ dual Verdiera
$\mathbb D_XA=R\mathcal Hom_X(A,\omega_X)$.
\end{definicja}
\begin{twierdzenie}[Dualność Verdiera dla ciała]
\label{thm:xvi-verdier}
Dla lokalnie zwartej, odpowiednio skończeniewymiarowej
przestrzeni $X$ i $A\in D^b_c(X)$ zachodzi naturalny
quasi-izomorfizm
\[
R\Gamma(X,\mathbb D_X A)
\simeq R\operatorname{Hom}_k(R\Gamma_c(X,A),k).
\]
Jeśli $X$ jest rzeczywistą rozmaitością wymiaru $n$,
$\omega_X\simeq\operatorname{or}_X[n]$, gdzie
$\operatorname{or}_X$ jest systemem orientacji.
\end{twierdzenie}
\begin{proof}[Dowód wskazanego modelu i wyraźna zależność]
Dla $X=\R^n$ i $A=k_X$ obie strony mają jedną
jednowymiarową grupę kohomologii w stopniu $-n$:
$R\Gamma_c(\R^n;k)\simeq k[-n]$, zatem
$R\operatorname{Hom}_k(k[-n],k)\simeq k[n]$.
Również $R\Gamma(\R^n,k_X[n])\simeq k[n]$.
Dla ogólnej rozmaitości te lokalne identyfikacje skleja
system orientacji. Istnienie kanonicznego quasi-izomorfizmu,
funktorialność i twierdzenie dla osobliwych $X$ należą do
pełnej dualności Verdiera; nie wynikają z porównania wymiarów
w przykładzie euklidesowym.
\end{proof}
\begin{wniosek}[Dualność Poincarégo jako szczególny przypadek]
Jeśli $M$ jest spójną, orientowalną rozmaitością rzeczywistą
wymiaru $n$, to dla ciała $k$ zachodzi
$H_c^q(M;k)^*\cong H^{n-q}(M;k)$ przy odpowiednim
sparowaniu dualności. Dla rozmaitości nieorientowalnej
prawa strona otrzymuje lokalne współczynniki orientacji.
\end{wniosek}
\begin{proof}
Wstawiamy $A=k_M$ do poprzedniego twierdzenia i
$\omega_M\simeq\operatorname{or}_M[n]$;
porównanie stopni kohomologicznych daje wynik.
\end{proof}
\begin{uwaga}[Uwaga o skończoności wymiarów]
Identyfikacja przez liniowy dual w powyższej postaci jest
naturalna nad ciałem, lecz nie oznacza automatycznie, że
wszystkie grupy kohomologii są skończeniewymiarowe.
W obliczeniach na skończonych kompleksach i na zwartych
odmianach algebraicznych problem ten nie występuje.
\end{uwaga}

\section{Dualność i obrazy z nośnikami właściwymi}
\begin{twierdzenie}[Wymiana obrazów przez dualność]
\label{thm:xvi-duality-exchange}
Niech $f:X\to Y$ będzie morfizmem w kategorii
konstruktybilnych kompleksów, a $A\in D^b_c(X)$.
Istnieje naturalny izomorfizm
\[
\mathbb D_Y(Rf_!A)\simeq Rf_*\mathbb D_X A.
\]
Jeśli $f$ jest właściwe, $Rf_!\simeq Rf_*$ i
$\mathbb D_Y Rf_*A\simeq Rf_*\mathbb D_XA$.
\end{twierdzenie}
\begin{proof}[Wyprowadzenie z adjunkcji]
Dla $B\in D^b_c(Y)$ rachunek formalny daje
\[
\begin{aligned}
R\operatorname{Hom}_Y(B,\mathbb D_Y Rf_!A)
&\simeq R\operatorname{Hom}_Y(B\otimes^L Rf_!A,\omega_Y)\\
&\simeq R\operatorname{Hom}_X(f^*B\otimes^L A,f^!\omega_Y)\\
&\simeq R\operatorname{Hom}_X(f^*B,\mathbb D_XA)\\
&\simeq R\operatorname{Hom}_Y(B,Rf_*\mathbb D_XA).
\end{aligned}
\]
W drugim kroku korzystamy ze wzoru projekcji i adjunkcji
$(Rf_!,f^!)$; w trzecim z kanonicznego
$f^!\omega_Y\simeq\omega_X$. Yoneda daje izomorfizm.
To formalny dowód \emph{warunkowy}: zależy od wcześniej
przyjętych twierdzeń o wzorze projekcji i istnieniu $f^!$.
\end{proof}

\section{Zmiana podstawy i wzór projekcji}
Rozważmy kartezjański kwadrat
\[
\begin{tikzpicture}[baseline=(current bounding box.center),node distance=22mm]
\node (Xp) {$X'$}; \node[right=of Xp] (X) {$X$};
\node[below=of Xp] (Yp) {$Y'$}; \node[right=of Yp] (Y) {$Y$};
\draw[->] (Xp)--node[above] {$g'$}(X);
\draw[->] (Xp)--node[left] {$f'$}(Yp);
\draw[->] (X)--node[right] {$f$}(Y);
\draw[->] (Yp)--node[below] {$g$}(Y);
\end{tikzpicture}
\]

\begin{twierdzenie}[Zmiana podstawy dla obrazu z nośnikami właściwymi]
\label{thm:xvi-basechange}
Przy standardowych założeniach formalizmu sześciu operacji
istnieje naturalny izomorfizm
$g^*Rf_!A\simeq Rf'_!(g')^*A$
dla konstruktybilnego $A$.
Jeśli $f$ jest właściwa, otrzymujemy też
$g^*Rf_*A\simeq Rf'_*(g')^*A$.
\end{twierdzenie}
\begin{proof}[Dowód w prostym przypadku i zakres twierdzenia]
Jeśli $f$ jest projekcją skończonego dyskretnego zbioru
$\{1,\dots,m\}\times Y\to Y$, to $f_!=f_*$ i
$(Rf_*A)_y\cong\bigoplus_{i=1}^m A_{(i,y)}$.
Identyfikacja pozostaje prawdziwa po dowolnym
$g:Y'\to Y$ i daje żądany izomorfizm na kiełkach.
Przypadek ogólny wymaga twierdzenia o zmianie podstawy
z właściwymi nośnikami; jest tu przyjęty, a nie
uzasadniony tym szczególnym obliczeniem.
\end{proof}

\begin{twierdzenie}[Wzór projekcji]
\label{thm:xvi-projection}
Jeżeli $f:X\to Y$ jest mapą rozważanej kategorii,
$A\in D^b_c(X)$ oraz $L$ jest lokalnym systemem
skończeniewymiarowych $k$-przestrzeni na $Y$, to
\[
Rf_!(A\otimes_k f^*L)\simeq (Rf_!A)\otimes_k L.
\]
Ogólniejszy wzór dla $B\in D^b_c(Y)$ ma postać
$Rf_!(A\otimes^L f^*B)\simeq Rf_!A\otimes^L B$,
przy zwykłych hipotezach konstruktybilności i
skończonego wymiaru kohomologicznego.
\end{twierdzenie}
\begin{proof}[Wersja dla lokalnego systemu]
Lokalnie na $Y$ mamy $L|_V\cong k_V^r$.
Dla $V$ otwartego po obu stronach dostajemy
odpowiednio $r$ kopii $Rf_!A|_V$, ponieważ
funktor $Rf_!$ jest addytywny, a tensor z $k^r$
odpowiada sumie skończonej. Izomorfizmy są naturalne
w zmianach trywializacji $L$, więc sklejają się na $Y$.
Uogólnienie na dowolne konstruktybilne $B$ nie wynika
z lokalnej trywialności $B$: jest odrębnym twierdzeniem
formalizmu sześciu operacji.
\end{proof}

\section{Włókna, obrazy pochodne i ciąg spektralny Leraya}
\begin{twierdzenie}[Kohomologia włókien przy mapie właściwej]
Niech $f:X\to Y$ będzie właściwą mapą rozważanej kategorii,
$A\in D^b_c(X;k)$, $y\in Y$, a $i_y:\{y\}\hookrightarrow Y$.
Dla każdego $q\in\Z$ istnieje naturalny izomorfizm
\[
\mathcal H^q(i_y^*Rf_*A)
\cong\mathbb H^q(f^{-1}(y);A|_{f^{-1}(y)}).
\]
W szczególności lewa strona jest kiełkiem
$\mathcal H^q(Rf_*A)$ w $y$. Dla snopu $\mathcal F$
umieszczonego w stopniu zerowym można ją zapisać
$(R^qf_*\mathcal F)_y$.
\end{twierdzenie}
\begin{proof}
Zastosujmy Twierdzenie~\ref{thm:xvi-basechange} do
$g:\{y\}\hookrightarrow Y$. Lewa strona jest
kompleksem $i_y^*Rf_*A$, a prawa
$R\Gamma(f^{-1}(y),A|_{f^{-1}(y)})$.
Po wzięciu kohomologii stopnia $q$ otrzymujemy tezę.
\end{proof}
\begin{uwaga}
Dla $f$ niewłaściwej analogiczny wzór z $Rf_*$ jest na ogół
fałszywy: wymaga dodatkowej hipotezy zmiany podstawy. Natomiast przy obrazach z nośnikami właściwymi
otrzymujemy lokalnie grupy $H_c^q(f^{-1}(y);A|_{f^{-1}(y)})$,
o ile spełniono założenia zmiany podstawy dla $Rf_!$.
\end{uwaga}
\begin{przyklad}[Złożenie $x\mapsto x^2$]
Dla właściwej mapy $f:\R\to\R$, $f(x)=x^2$, włókno
nad $y<0$ jest puste, nad $y=0$ ma jeden punkt, a nad
$y>0$ dwa punkty. Zatem snop $R^0f_*k_\R$ ma
odpowiednio rzędy $0,1,2$, a $R^qf_*k_\R=0$ dla $q>0$.
To konstruktybilny snop, który nie jest lokalnym systemem
na całej prostej.
\end{przyklad}
\begin{twierdzenie}[Ciąg Leraya]
Dla mapy $f:X\to Y$ w rozważanej kategorii i snopu
$\mathcal F$ istnieje pierwszoćwiartkowy ciąg spektralny
$E_2^{p,q}=H^p(Y,R^qf_*\mathcal F)
\Longrightarrow H^{p+q}(X,\mathcal F)$.
Dla $Rf_!$ i $R\Gamma_c$ istnieje odpowiedni ciąg
spektralny z nośnikami zwartymi, przy standardowych
założeniach skończoności.
\end{twierdzenie}
\begin{proof}[Wyjaśnienie konstrukcji]
Złożenie lewostronnie dokładnych funktorów
$\Gamma(Y,-)\circ f_*=\Gamma(X,-)$ prowadzi do
ciągu Grothendiecka. Dla snopów injektywnych
$f_*$ jest acykliczny względem $\Gamma(Y,-)$:
$\operatorname{Hom}(-,f_*I)\cong
\operatorname{Hom}(f^*(-),I)$ i dokładność $f^*$
pokazują, że $f_*I$ jest injektywny.
Reszta konstrukcji ciągu Grothendiecka to wynik
z algebry homologicznej rozwinięty wcześniej.
\end{proof}

\section{Cykle pobliskie i zanikające: konstrukcja z monodromią}
\label{sec:xvi-nearby}
Niech $f:X\to\Delta$ będzie holomorficzną funkcją,
$X_0=f^{-1}(0)$, $X^*=f^{-1}(\Delta^*)$, a
$i:X_0\hookrightarrow X$ oraz $j:X^*\hookrightarrow X$
odpowiednimi włożeniami. Oznaczmy przez
$\pi:\widetilde{\Delta^*}\to\Delta^*$ uniwersalne nakrycie
i przez $\widetilde X^*=X^*\times_{\Delta^*}
\widetilde{\Delta^*}$ odpowiednią przestrzeń. Niech
$\widetilde j:\widetilde X^*\to X$ będzie naturalną mapą.

\begin{definicja}[Cykle pobliskie i zanikające]
Dla $A\in D^b_c(X)$ określamy
$R\psi_fA=i^*R\widetilde j_*\widetilde j^*A$
na $X_0$; w tym zapisie $\widetilde j^*A$ oznacza
pullback $A$ na $\widetilde X^*$.
Morfizm specjalizacji $\mathrm{sp}:i^*A\to R\psi_fA$
wyznacza kompleks cykli zanikających
$R\phi_fA=\operatorname{Cone}(\mathrm{sp})$.
Otrzymujemy trójkąt
\[
 i^*A\xrightarrow{\mathrm{sp}}R\psi_fA
 \longrightarrow R\phi_fA\longrightarrow i^*A[1].
\]
Deck transformation nakrycia indukuje operator monodromii
na $R\psi_fA$ i $R\phi_fA$.
\end{definicja}
\begin{uwaga}[Znaczenie konwencji]
W niektórych książkach używa się przesuniętego
kompleksu $R\phi_f[-1]$ lub normalizacji perwersyjnej.
W naszym tekście pozostajemy przy stożku bez przesunięcia,
spójnym z Rozdziałem~\ref{ch:xiv-hodge-monodromia}.
\end{uwaga}
\begin{stwierdzenie}[Włókno Milnora jako kiełek cykli pobliskich]
\label{prop:xvi-milnor-stalk}
Jeśli $x\in X_0$ jest punktem, w którym $f$
ma izolowaną osobliwość krytyczną, to dla małej kuli
$B_\varepsilon(x)$ i dostatecznie małego $t\ne0$
kiełek $(R^q\psi_fk_X)_x$ można zidentyfikować
z $H^q(B_\varepsilon(x)\cap f^{-1}(t);k)$.
\end{stwierdzenie}
\begin{proof}[Status]
Wybór promieni i parametrów jest taki jak w
Twierdzeniu Milnora: najpierw ustala się
$\varepsilon$, a następnie odpowiednio małe $|t|$.
Identyfikacja wykorzystuje lokalną trywialność
fibracji Milnora i odpowiednie porównanie
pochodnego obrazu na uniwersalnym nakryciu.
Pełny dowód tego twierdzenia porównawczego byłby
powtórzeniem znacznej części teorii Milnora;
przyjmujemy go jawnie jako twierdzenie pomocnicze.
\end{proof}
\begin{przyklad}[Węzeł $zw=t$]
Dla $f(z,w)=zw$ włókno Milnora w małej kuli jest
homotopijnie równoważne $S^1$, zatem
$(R^0\psi_fk_X)_0\cong k$ i
$(R^1\psi_fk_X)_0\cong k$, a pozostałe grupy znikają.
Morfizm $i^*k_X\to R\psi_fk_X$ jest izomorfizmem
w stopniu zerowym, więc
$(R^1\phi_fk_X)_0\cong k$, pozostałe grupy znikają.
Lokalna monodromia na $H^1(S^1;k)$ jest tożsamością,
co nie przeczy nietrywialnej transwekcji na
$H^1$ globalnego gładkiego włókna zwartej rodziny.
\end{przyklad}
\begin{przyklad}[Rozgałęzienie $z^2=t$]
Dla $f:\C\to\C$, $f(z)=z^2$, włókno Milnora w
punkcie $0$ składa się z dwóch punktów.
Otrzymujemy $(R^0\psi_fk_\C)_0\cong k^2$, a
morfizm specjalizacji w stopniu zerowym jest przekątną
$k\to k^2$. Zatem $(R^0\phi_fk_\C)_0\cong k$.
Po okrążeniu $t=0$ dwa punkty włókna zamieniają się,
więc na zredukowanej $H^0$ monodromia działa przez $-1$
(gdy $\operatorname{char}k=2$, znak ten jest równy $+1$).
\end{przyklad}

\section{Lokalny schemat zależności i ograniczenia formalizmu}
\begin{center}
\begin{tikzpicture}[node distance=12mm and 13mm,>=Stealth,
 box/.style={draw,rounded corners,align=center,inner sep=5pt,font=\small}]
\node[box] (a) {$A\in D^b_c(X)$};
\node[box,below left=of a] (b) {$Rf_*,\ Rf_!$};
\node[box,below right=of a] (c) {$f^*,\ f^!$};
\node[box,below=25mm of a] (d) {$\mathbb D_X$ i $\omega_X$};
\node[box,below=of d] (e) {lokalizacja, zmiana podstawy\\cykle $R\psi_f,\ R\phi_f$};
\draw[->] (a)--(b); \draw[->] (a)--(c);
\draw[->] (b)--(d); \draw[->] (c)--(d); \draw[->] (d)--(e);
\end{tikzpicture}
\par\smallskip\small Powiązania sześciu operacji. Tensor i wewnętrzny Hom
są obecne w formule projekcji i definicji dualności.
\end{center}

\begin{uwaga}[Topologia klasyczna a topologia \emph{étale}]
W tym rozdziale pracujemy ze snopami na klasycznej przestrzeni
analitycznej nad $\C$. W geometrii algebraicznej podobne formuły
występują dla snopów \emph{étale}, lecz ich dowody i ograniczenia
(w szczególności warunki na torsję i charakterystykę) są inne.
Samo podobieństwo symboli $f^!,Rf_!$ nie upoważnia do
przeniesienia twierdzenia bez sprawdzenia kategorii współczynników.
\end{uwaga}
\begin{uwaga}[Precyzja stosowania funktorów]
Dla dowolnego kompleksu $A$ wyrażenie $R^qf_*A$ jest
skrótowe i bywa niejednoznaczne; używamy
$\mathcal H^q(Rf_*A)$. Izomorfizmy w rodzaju
$g^*Rf_!\simeq Rf'_!(g')^*$ lub
$\mathbb D_YRf_!\simeq Rf_*\mathbb D_X$ odnosimy do
kategorii i hipotez wskazanych przed twierdzeniami,
nie do wszystkich odwzorowań między dowolnymi
przestrzeniami topologicznymi. Przy pracy ze snopami
koherentnymi na schematach te same symbole mogą oznaczać
inne operacje, zwłaszcza pochodny pullback tensorowy.
\end{uwaga}
\begin{uwaga}[Status dowodów]
W rozdziale samodzielnie uzasadniono długi ciąg stożka,
krótki ciąg lokalizacyjny, podstawową adjunkcję $f^*\dashv Rf_*$,
przypadek punktu na prostej, wersję projekcji dla
lokalnych systemów oraz obliczenia dla $x^2$, $z^2$ i $zw$.
Istnienie ogólnego $f^!$, twierdzenie o zmianie podstawy,
pełna dualność Verdiera, ogólny wzór projekcji i formalne
własności cykli pobliskich wymagają niezależnych wyników
formalizmu sześciu operacji. Zostały przyjęte jawnie,
nie jako konsekwencje elementarnych rachunków.
\end{uwaga}

